% part: intuitionistic-logic
% chapter: semantics
% section: introduction

\documentclass[../../../include/open-logic-section]{subfiles}

\begin{document}

\olfileid{int}{sem}{int}

\olsection{مقدمة}

لا يوصف أي منطق وصفًا مُرضيًا من دون دلالات، والمنطق الحدسي ليس استثناءً.
وفي حين تُعد الدلالات القائمة على ال!!{valuation}s قانونية للمنطق
الكلاسيكي، توجد للمنطق الحدسي دلالات متنافسة عدة. ولا تُرضي أي منها إرضاءً
تامًا بمعنى أنها تقدم تفسيرًا لمعاني الروابط مقبولًا من وجهة النظر الحدسية.

ولعل الدلالات القائمة على ال!!{relational model}s، المشابهة لدلالات المنطق
الجهوي، أشهرها. ففي هذه الدلالات تُسند ال!!{propositional variable}s إلى
عوالم، وترتبط هذه العوالم بعلاقة نفاذ. وتكون هذه العلاقة دائمًا ترتيبًا
جزئيًا؛ أي إنها انعكاسية ومضادة للتناظر ومتعدية.

ويمكنك حدسيًا أن تتصور هذه العوالم حالات للمعرفة أو «أوضاعًا استدلالية».
وتكون الحالة~$w'$ متاحة من~$w$ إذا وفقط إذا كانت، في حدود علمنا، حالة
معرفية (مستقبلية) ممكنة، أي متوافقة مع ما هو معلوم عند~$w$. ومتى عُرفت
قضية، لم تعد قابلة لأن تصبح غير معلومة؛ أي كلما كانت~$!A$ معلومة عند~$w$
وكان~$Rww'$، كانت~$!A$ معلومة عند~$w'$ أيضًا. ومن ثم تكون «المعرفة»
رتيبة بالنسبة إلى علاقة النفاذ.

إذا عرّفنا «كون~$!A$ معلومة»، كما في منطق المعرفة، بأنه «صدقها في جميع
البدائل المعرفية»، كانت~$!A \land !B$ معلومة عند~$w$ إذا كانت~$!A$
و$!B$ معلومتين معًا في جميع البدائل المعرفية. لكن لما كانت المعرفة رتيبة
وكانت~$R$ انعكاسية، كان معنى ذلك أن~$!A \land !B$ معلومة عند~$w$ إذا
وفقط إذا كانت~$!A$ و$!B$ معلومتين عند~$w$. وللسبب نفسه، تكون
$!A \lor !B$ معلومة عند~$w$ إذا وفقط إذا كانت إحداهما على الأقل معلومة.
ومن ثم تتطابق شروط صدق الرابطين~$\land$ و$\lor$ مع شروطهما في المنطق
الكلاسيكي.

أما شروط صدق الشرطية فتختلف عن المنطق الكلاسيكي. تكون~$!A \lif !B$
معلومة عند~$w$ إذا وفقط إذا لم توجد~$w'$ بحيث~$Rww'$ تكون عندها~$!A$
معلومة من دون أن تكون~$!B$ معلومة أيضًا. وهذا لا يساوي شرط كون~$!A$ غير
معلومة أو~$!B$ معلومة عند~$w$. فإذا لم نكن نعلم~$!A$ ولا~$!B$ عند~$w$،
فقد توجد حالة معرفية مستقبلية~$w'$ بحيث~$Rww'$، تصبح عندها~$!A$ معلومة
من دون أن تصبح~$!B$ معلومة أيضًا.

ولا نعلم~$\lnot !A$ إلا إذا لم توجد حالة معرفية مستقبلية ممكنة نعلم فيها
$!A$. والفكرة هنا أنه لو أمكن أن تُعلم~$!A$، لصارت~$!A$ معلومة في حالة
معرفية مستقبلية ممكنة ما. ولأننا لا نستطيع معرفة~$\lfalse$، لكنا في تلك
الحالة المعرفية المستقبلية نعلم~$!A$ ولا نعلم~$\lfalse$.

وعلى هذا التفسير يفشل مبدأ الثالث المرفوع. فثمة~$!A$ لا نعلمها بعد، لكن قد
نعلمها مستقبلًا. وبالنسبة إلى !!a{formula} كهذه~$!A$، تكون~$!A$ و
$\lnot !A$ كلتاهما غير معلومة، ولذلك لا تكون~$!A \lor \lnot !A$
معلومة. لكننا نعلم، مثلًا، أن~$\lnot(!A \land \lnot !A)$. إذ لا يمكن أن
توجد حالة مستقبلية نعلم فيها~$!A$ و$\lnot !A$ معًا، ونعلم ذلك بصرف النظر
عن علمنا أو عدم علمنا بـ~$!A$ أو~$\lnot !A$.

ليست ال!!^{relational model}s الدلالات الوحيدة المتاحة للمنطق الحدسي؛
فالدلالات الطوبولوجية بديل آخر. وفيها تُفسر القضايا بوصفها مجموعات مفتوحة
في فضاء طوبولوجي، وتُفسر الروابط بوصفها عمليات على هذه المجموعات (فمثلًا
يقابل~$\land$ التقاطع).

\end{document}
